\documentclass[11pt,a4paper]{article}
% preamble.tex — shared setup for human-verification documents
% Compile each document with:  tectonic <File>.tex

\usepackage[T1]{fontenc}
\usepackage[utf8]{inputenc}
\usepackage{lmodern}
\usepackage{amssymb,amsmath,amsthm}
\usepackage{mathtools}
\usepackage[dvipsnames]{xcolor}
\usepackage{listings}
\usepackage{hyperref}
\usepackage[margin=2.5cm]{geometry}
\usepackage{parskip}

% ── Colours ──────────────────────────────────────────────────────────
\definecolor{leanBg}{HTML}{F7F7F0}
\definecolor{haskellBg}{HTML}{F0F4F7}
\definecolor{keyword}{HTML}{7B3294}
\definecolor{comment}{HTML}{5A8C32}
\definecolor{noteFg}{HTML}{B35900}

% ── Lean 4 listings language ─────────────────────────────────────────
% NOTE: Lean listing blocks use ASCII approximations of unicode symbols
% (e.g. -> for arrow, o. for compose).  The literate key substitutes
% these with proper math-mode renderings in the PDF.
\lstdefinelanguage{Lean4}{
  morekeywords={def,theorem,structure,where,namespace,end,import,fun,
                let,have,show,by,exact,intro,cases,rfl,subst,obtain,
                funext,simp,rw,apply,variable,Type,Prop,instance,
                section,open,class,extends,deriving,sorry,noncomputable},
  sensitive=true,
  morecomment=[l]{--},
  morecomment=[n]{/-}{-/},
  morestring=[b]",
}

% ── Lean listing style ───────────────────────────────────────────────
\lstdefinestyle{lean}{
  language=Lean4,
  backgroundcolor=\color{leanBg},
  basicstyle=\small\ttfamily,
  keywordstyle=\color{keyword}\bfseries,
  commentstyle=\color{comment}\itshape,
  breaklines=true,
  frame=single,
  framerule=0.4pt,
  xleftmargin=4pt,
  xrightmargin=4pt,
  aboveskip=6pt,
  belowskip=4pt,
  literate=
    {->}{{$\to$}}2
    {=>}{{$\Rightarrow$}}2
    {<|}{{$\langle$}}1
    {|>}{{$\rangle$}}1
    {==>}{{$\Longrightarrow$}}3
    { o. }{{$\;\circ\;$}}3
    {forall}{{$\forall$}}1
    {:=}{{$\,:=\,$}}2
    {!=}{{$\neq$}}2
    {<=}{{$\leq$}}2
    {>=}{{$\geq$}}2,
}

% ── Haskell listing style ───────────────────────────────────────────
\lstdefinestyle{haskell}{
  language=Haskell,
  backgroundcolor=\color{haskellBg},
  basicstyle=\small\ttfamily,
  keywordstyle=\color{keyword}\bfseries,
  commentstyle=\color{comment}\itshape,
  breaklines=true,
  frame=single,
  framerule=0.4pt,
  xleftmargin=4pt,
  xrightmargin=4pt,
  aboveskip=6pt,
  belowskip=4pt,
  literate=
    {->}{{$\to$}}2
    {=>}{{$\Rightarrow$}}2,
}

% ── Theorem environments ────────────────────────────────────────────
\theoremstyle{definition}
\newtheorem{defn}{Definition}[section]
\newtheorem{thm}[defn]{Theorem}
\newtheorem{prop}[defn]{Proposition}

\theoremstyle{remark}
\newtheorem*{fidelitynote}{Fidelity note}

% ── Convenience macros ───────────────────────────────────────────────
\newcommand{\lean}[1]{\texttt{#1}}      % inline Lean identifier
\newcommand{\hask}[1]{\texttt{#1}}      % inline Haskell identifier
\newcommand{\Ext}{\mathrm{Ext}}
\newcommand{\Pos}{\mathrm{Pos}}
\newcommand{\Shape}{\mathrm{Shape}}


\title{Human Verification: \texttt{Containers.DContCat}\\[4pt]
  \large Directed Containers as Categories, and Cofunctors}
\author{Machine-generated from \texttt{lean/DContCat.lean} for review by Neil Ghani}
\date{July 2026}

\begin{document}
\maketitle

\begin{abstract}
This document presents every definition and theorem \emph{statement} from the
Lean~4 module \lean{Containers.DContCat} in three parallel forms: the original
Lean, a readable Haskell rendering, and precise mathematical English.  Proofs
are omitted --- they are machine-checked by Lean.  The purpose is to let a
human verifier confirm that the formalisation captures the intended
mathematics.
\end{abstract}

\tableofcontents
\bigskip

%===================================================================
\section{Overview}
%===================================================================

This file formalises the \emph{object-level dictionary} of Ahman--Uustalu
(2016): a directed container $D$ presents a small category $\mathcal{C}(D)$,
and conversely every small category presents a directed container.  It also
formalises the \emph{morphism-level} result: directed-container morphisms are
exactly \textbf{cofunctors} between the induced categories.

The key points:
\begin{itemize}
  \item Objects of $\mathcal{C}(D)$ are shapes; morphisms $s \to s'$ are
        positions $p : P(s)$ with $s \downarrow p = s'$.
  \item The identity at $s$ is the root $o(s)$; composition is the shift
        $\oplus$.
  \item Category laws are exactly D2/D3 (unit laws) and D5 (associativity).
  \item The bijection between directed-container morphisms and cofunctors is
        witnessed by \lean{rfl} in both directions --- the two structures carry
        literally the same data under different field names.
\end{itemize}

%===================================================================
\section{Small categories}
%===================================================================

% --- 1. SmallCat ---------------------------------------------------
\subsection{Structure: \lean{SmallCat}}

\begin{defn}[Small category]
A \textbf{small category} $\mathcal{C}$ consists of:
\begin{itemize}
  \item a type $\mathrm{Obj}$ of \emph{objects};
  \item for each pair $a, b : \mathrm{Obj}$, a type $\mathrm{Hom}(a, b)$ of
        \emph{morphisms};
  \item for each $a$, an \emph{identity} $\mathrm{id}_a : \mathrm{Hom}(a, a)$;
  \item \emph{diagrammatic composition}
        $\mathrm{comp} : \mathrm{Hom}(a,b) \to \mathrm{Hom}(b,c) \to
        \mathrm{Hom}(a,c)$
        (left-to-right order, matching $\oplus$ in directed containers);
\end{itemize}
satisfying:
\begin{itemize}
  \item \textbf{left unit:} $\mathrm{comp}(\mathrm{id}_a, f) = f$;
  \item \textbf{right unit:} $\mathrm{comp}(f, \mathrm{id}_b) = f$;
  \item \textbf{associativity:}
        $\mathrm{comp}(\mathrm{comp}(f,g), h) = \mathrm{comp}(f, \mathrm{comp}(g,h))$.
\end{itemize}
\end{defn}

\paragraph{Lean.}
\begin{lstlisting}[style=lean]
structure SmallCat where
  Obj : Type
  Hom : Obj -> Obj -> Type
  id : (a : Obj) -> Hom a a
  comp : {a b c : Obj} -> Hom a b -> Hom b c -> Hom a c
  id_comp : forall {a b : Obj} (f : Hom a b), comp (id a) f = f
  comp_id : forall {a b : Obj} (f : Hom a b), comp f (id b) = f
  assoc : forall {a b c d : Obj} (f : Hom a b) (g : Hom b c) (h : Hom c d),
    comp (comp f g) h = comp f (comp g h)
\end{lstlisting}

\paragraph{Haskell.}
\begin{lstlisting}[style=haskell]
-- A small category, hand-rolled.
-- Composition is in diagrammatic (left-to-right) order.
data SmallCat obj = SmallCat
  { hom     :: obj -> obj -> Type  -- not expressible in value-level Haskell
  , catId   :: obj -> ()           -- identity (approximated)
  , catComp :: ()                  -- composition (approximated)
  }
-- The laws are stated as properties (QuickCheck style):
-- prop_id_comp  f = comp (catId a) f == f
-- prop_comp_id  f = comp f (catId b) == f
-- prop_assoc f g h = comp (comp f g) h == comp f (comp g h)
\end{lstlisting}

\begin{fidelitynote}
Lean's structure bundles the proof fields directly.  Haskell cannot express
dependent hom-types or carry proof fields in a data type without significant
ceremony; the Haskell rendering approximates the data and states the laws as
external properties.
\end{fidelitynote}


%===================================================================
\section{The dictionary: directed container \texorpdfstring{$\leftrightarrow$}{<->} small category}
%===================================================================

% --- 2. Helper lemma: transport_cancel ----------------------------
\subsection{Lemma: \lean{transport\_cancel}}

\begin{prop}[Transport round-trip]
For any type family $P : \alpha \to \mathbf{Type}$ and equality $h : x = y$,
transporting a term $p : P(y)$ back along $h$ and then forward along $h$ is
the identity:
\[
  h \triangleright h^{-1} \triangleright p \;=\; p.
\]
\end{prop}

\paragraph{Lean.}
\begin{lstlisting}[style=lean]
theorem transport_cancel {alpha : Type} {P : alpha -> Type} {x y : alpha}
    (h : x = y) (p : P y) : h [> h.symm [> p = p
\end{lstlisting}

\paragraph{Haskell.}
\begin{lstlisting}[style=haskell]
-- Transport is trivial in Haskell (no dependent types).
-- The round-trip is the identity by definition.
\end{lstlisting}


% --- 3. Helper lemma: sub_shift_eq --------------------------------
\subsection{Lemma: \lean{DirectedContainer.sub\_shift\_eq}}

\begin{prop}[Codomain of a composite morphism]
Let $D$ be a directed container.  If $p : P(a)$ with $a \downarrow p = b$,
and $q : P(b)$ with $b \downarrow q = c$, then
\[
  a \downarrow (p \oplus q') \;=\; c,
\]
where $q' = h_1^{-1} \triangleright q$ is $q$ transported to live over $a \downarrow p$.
\end{prop}

\paragraph{Lean.}
\begin{lstlisting}[style=lean]
theorem sub_shift_eq {a b c : D.Shape} (p : D.Pos a) (hp : D.sub a p = b)
    (q : D.Pos b) (hq : D.sub b q = c) :
    D.sub a (D.shift a p (hp.symm [> q)) = c
\end{lstlisting}

\paragraph{Haskell.}
\begin{lstlisting}[style=haskell]
-- Follows immediately from D4 when hp and hq are rfl.
-- sub_shift_eq :: sub a (shift a p q) == c
-- given: sub a p = b, sub b q = c
\end{lstlisting}


% --- 4. toSmallCat ------------------------------------------------
\subsection{Definition: \lean{DirectedContainer.toSmallCat}}

\begin{defn}[Small category presented by a directed container]
Given a directed container $D = (S, P, o, {\downarrow}, {\oplus})$, define a
small category $\mathcal{C}(D)$ as follows:
\begin{align*}
  \mathrm{Obj} &= S, \\
  \mathrm{Hom}(a, b) &= \{ p : P(a) \mid a \downarrow p = b \}, \\
  \mathrm{id}_a &= (o(a),\; D_1(a)), \\
  \mathrm{comp}(p, q) &= (p \oplus q',\; \text{\lean{sub\_shift\_eq}}),
\end{align*}
where $q'$ is $q$ transported from the fibre over $b$ to the fibre over
$a \downarrow p$ via the equality $p.\mathrm{property}$.
The unit and associativity laws are D2/D3 and D5 respectively (after
substituting the codomain equalities).
\end{defn}

\paragraph{Lean.}
\begin{lstlisting}[style=lean]
def toSmallCat : SmallCat where
  Obj := D.Shape
  Hom a b := { p : D.Pos a // D.sub a p = b }
  id a := <|D.root a, D.d1 a|>
  comp f g :=
    <|D.shift _ f.val (f.property.symm [> g.val),
     D.sub_shift_eq f.val f.property g.val g.property|>
  id_comp := ...
  comp_id := ...
  assoc := ...
\end{lstlisting}

\paragraph{Haskell.}
\begin{lstlisting}[style=haskell]
-- The category presented by a directed container.
-- Objects: shapes s
-- Morphisms s->t: positions p over s with sub s p = t
-- Identity at s: root s (with D1 saying sub s (root s) = s)
-- Composition (p, q) where sub s p = t: shift s p q
--   (q transported to live over sub s p)
\end{lstlisting}


% --- 5. SmallCat.toDirectedContainer ------------------------------
\subsection{Definition: \lean{SmallCat.toDirectedContainer}}

\begin{defn}[Directed container presented by a small category]
Given a small category $\mathcal{C}$, define a directed container as follows:
\begin{align*}
  S &= \mathrm{Obj}, \\
  P(s) &= \Sigma_{s'}\, \mathcal{C}(s, s'), \\
  o(s) &= (s, \mathrm{id}_s), \\
  s \downarrow p &= p.1 \quad\text{(codomain projection)}, \\
  p \oplus q &= (q.1,\; \mathrm{comp}(p.2, q.2)).
\end{align*}
Laws D1 and D4 hold by $\mathrm{rfl}$ (the codomain of a composite is the
codomain of the second factor, definitionally).  D2/D3/D5 are the category
laws.
\end{defn}

\paragraph{Lean.}
\begin{lstlisting}[style=lean]
def toDirectedContainer (C : SmallCat) : DirectedContainer where
  Shape := C.Obj
  Pos s := (s' : C.Obj) * C.Hom s s'
  root s := <|s, C.id s|>
  sub _ p := p.1
  shift _ p q := <|q.1, C.comp p.2 q.2|>
  d1 _ := rfl
  d4 _ _ _ := rfl
  d2 _ q := congrArg (Sigma.mk q.1) (C.id_comp q.2)
  d3 _ p := congrArg (Sigma.mk p.1) (C.comp_id p.2)
  d5 _ p q r := congrArg (Sigma.mk r.1) (C.assoc p.2 q.2 r.2)
\end{lstlisting}

\paragraph{Haskell.}
\begin{lstlisting}[style=haskell]
-- Directed container from a small category.
data DCont obj = DCont
  { shape  :: obj
  , pos    :: obj -> [(obj, Morphism)]  -- outgoing morphisms
  , root   :: obj -> (obj, Identity)
  , sub    :: obj -> (obj, Morphism) -> obj   -- = fst (codomain)
  , shift  :: obj -> (obj, Morphism) -> (obj, Morphism) -> (obj, Morphism)
  }
-- shift _ (s', f) (s'', g) = (s'', comp f g)
\end{lstlisting}


%===================================================================
\section{The round-trip, honestly}
%===================================================================

The two constructions are \emph{not} definitionally mutually inverse.  Going
$D \mapsto \mathcal{C}(D) \mapsto D'$, the positions of $D'$ at $s$ are
$\Sigma_{s'}\, \{p : P(s) \mid s \downarrow p = s'\}$ --- each position tagged
with its codomain --- whereas $D$ has plain $P(s)$.  The two are in bijection
(a position determines its codomain), exhibited below.

% --- 6. toRoundtripPos / ofRoundtripPos ---------------------------
\subsection{Definitions: \lean{toRoundtripPos} and \lean{ofRoundtripPos}}

\begin{defn}[Round-trip position bijection (forward)]
Tag a position with its codomain:
\[
  \mathrm{toRoundtripPos}(s, p) \;=\; (s \downarrow p,\; p,\; \mathrm{rfl}).
\]
\end{defn}

\begin{defn}[Round-trip position bijection (backward)]
Forget the codomain tag:
\[
  \mathrm{ofRoundtripPos}(s, (s', p, h)) \;=\; p.
\]
\end{defn}

\paragraph{Lean.}
\begin{lstlisting}[style=lean]
def toRoundtripPos (D : DirectedContainer) (s : D.Shape) (p : D.Pos s) :
    D.toSmallCat.toDirectedContainer.Pos s :=
  <|D.sub s p, p, rfl|>

def ofRoundtripPos (D : DirectedContainer) (s : D.Shape)
    (p : D.toSmallCat.toDirectedContainer.Pos s) : D.Pos s :=
  p.2.val
\end{lstlisting}

\paragraph{Haskell.}
\begin{lstlisting}[style=haskell]
toRoundtripPos :: Pos s -> (Sub s, Pos s, ())   -- tag with codomain
toRoundtripPos p = (sub p, p, ())

ofRoundtripPos :: (Sub s, Pos s, ()) -> Pos s   -- forget tag
ofRoundtripPos (_, p, _) = p
\end{lstlisting}


% --- 7. Round-trip theorems ---------------------------------------
\subsection{Theorems: round-trip bijection}

\begin{thm}[Forward-backward is identity]
\[
  \mathrm{ofRoundtripPos}(s,\; \mathrm{toRoundtripPos}(s, p)) \;=\; p.
\]
\end{thm}

\begin{thm}[Backward-forward is identity]
\[
  \mathrm{toRoundtripPos}(s,\; \mathrm{ofRoundtripPos}(s, p)) \;=\; p.
\]
\end{thm}

\paragraph{Lean.}
\begin{lstlisting}[style=lean]
theorem ofRoundtripPos_toRoundtripPos (D : DirectedContainer)
    (s : D.Shape) (p : D.Pos s) :
    D.ofRoundtripPos s (D.toRoundtripPos s p) = p

theorem toRoundtripPos_ofRoundtripPos (D : DirectedContainer)
    (s : D.Shape) (p : D.toSmallCat.toDirectedContainer.Pos s) :
    D.toRoundtripPos s (D.ofRoundtripPos s p) = p
\end{lstlisting}

\paragraph{Haskell.}
\begin{lstlisting}[style=haskell]
-- prop_fwd_bwd p = ofRoundtripPos (toRoundtripPos p) == p
-- prop_bwd_fwd p = toRoundtripPos (ofRoundtripPos p) == p
\end{lstlisting}


%===================================================================
\section{Morphisms: directed-container morphisms and cofunctors}
%===================================================================

% --- 8. DContMorphism ---------------------------------------------
\subsection{Structure: \lean{DContMorphism}}

\begin{defn}[Directed-container morphism]
A \textbf{morphism of directed containers} $\varphi : C \Rightarrow D$ consists
of:
\begin{itemize}
  \item a \emph{covariant} shape map $f : S_C \to S_D$;
  \item a \emph{contravariant} position map $f^\sharp_s : P_D(f(s)) \to P_C(s)$;
\end{itemize}
satisfying:
\begin{itemize}
  \item \textbf{M0:} $f(C \!\downarrow_s\! f^\sharp_s(h)) = D \!\downarrow_{f(s)}\! h$;
  \item \textbf{M1:} $f^\sharp_s(o_D(f(s))) = o_C(s)$;
  \item \textbf{M2:} $f^\sharp_s(h \oplus_D k) = f^\sharp_s(h) \oplus_C f^\sharp_{C \downarrow_s f^\sharp_s(h)}(M0^{-1} \triangleright k)$.
\end{itemize}
\end{defn}

\paragraph{Lean.}
\begin{lstlisting}[style=lean]
structure DContMorphism (C D : DirectedContainer) where
  onShapes : C.Shape -> D.Shape
  onPos : (s : C.Shape) -> D.Pos (onShapes s) -> C.Pos s
  m0 : forall (s : C.Shape) (h : D.Pos (onShapes s)),
    onShapes (C.sub s (onPos s h)) = D.sub (onShapes s) h
  m1 : forall s : C.Shape, onPos s (D.root (onShapes s)) = C.root s
  m2 : forall (s : C.Shape) (h : D.Pos (onShapes s))
    (k : D.Pos (D.sub (onShapes s) h)),
    onPos s (D.shift (onShapes s) h k)
      = C.shift s (onPos s h)
          (onPos (C.sub s (onPos s h)) ((m0 s h).symm [> k))
\end{lstlisting}

\paragraph{Haskell.}
\begin{lstlisting}[style=haskell]
data DContMorph c d = DContMorph
  { onShapes :: Shape c -> Shape d      -- covariant
  , onPos    :: Shape c -> Pos d -> Pos c   -- contravariant
  -- M0: onShapes (sub_C s (onPos s h)) = sub_D (onShapes s) h
  -- M1: onPos s (root_D (onShapes s)) = root_C s
  -- M2: onPos s (shift_D (onShapes s) h k)
  --       = shift_C s (onPos s h)
  --           (onPos (sub_C s (onPos s h)) k)
  }
\end{lstlisting}

\begin{fidelitynote}
M2 in Lean involves a transport: $k : P_D(D \!\downarrow\! f(s) \!\downarrow\! h)$
must be transported along M0 to live over $f(C \!\downarrow\! f^\sharp_s(h))$ before
it can be fed to $\oplus_D$.  This transport is invisible in the Haskell
rendering since Haskell does not track which shape a position belongs to
at the type level.
\end{fidelitynote}


% --- 9. Cofunctor -------------------------------------------------
\subsection{Structure: \lean{Cofunctor}}

\begin{defn}[Cofunctor]
A \textbf{cofunctor} $\varphi : \mathcal{C}(C) \to \mathcal{C}(D)$ consists of:
\begin{itemize}
  \item an \emph{object map} $\mathrm{obj} : S_C \to S_D$;
  \item a \emph{lift} $\mathrm{lift}_s : P_D(\mathrm{obj}(s)) \to P_C(s)$
        (morphisms lift \emph{backwards});
\end{itemize}
satisfying the laws C0/C1/C2, which are word-for-word M0/M1/M2 of a directed-container
morphism.  The two structures carry \emph{literally the same data}.
\end{defn}

\paragraph{Lean.}
\begin{lstlisting}[style=lean]
structure Cofunctor (C D : DirectedContainer) where
  obj : C.Shape -> D.Shape
  lift : (s : C.Shape) -> D.Pos (obj s) -> C.Pos s
  c0 : forall (s : C.Shape) (h : D.Pos (obj s)),
    obj (C.sub s (lift s h)) = D.sub (obj s) h
  c1 : forall s : C.Shape, lift s (D.root (obj s)) = C.root s
  c2 : forall (s : C.Shape) (h : D.Pos (obj s))
    (k : D.Pos (D.sub (obj s) h)),
    lift s (D.shift (obj s) h k)
      = C.shift s (lift s h)
          (lift (C.sub s (lift s h)) ((c0 s h).symm [> k))
\end{lstlisting}

\paragraph{Haskell.}
\begin{lstlisting}[style=haskell]
-- Literally the same data as DContMorph, different field names.
data Cofunctor c d = Cofunctor
  { obj  :: Shape c -> Shape d
  , lift :: Shape c -> Pos d -> Pos c
  }
\end{lstlisting}


% --- 10. The hom-set bijection -----------------------------------
\subsection{Definitions and theorems: the hom-set bijection}

\begin{defn}[From morphism to cofunctor]
\lean{DContMorphism.toCofunctor}: the identity on data (M0/M1/M2 become
C0/C1/C2 by renaming).
\end{defn}

\begin{defn}[From cofunctor to morphism]
\lean{Cofunctor.toDContMorphism}: the identity on data (C0/C1/C2 become
M0/M1/M2 by renaming).
\end{defn}

\begin{thm}[Bijection, one way]
$\varphi.\lean{toCofunctor.toDContMorphism} = \varphi$ (by \lean{rfl}).
\end{thm}

\begin{thm}[Bijection, other way]
$\varphi.\lean{toDContMorphism.toCofunctor} = \varphi$ (by \lean{rfl}).
\end{thm}

\paragraph{Lean.}
\begin{lstlisting}[style=lean]
def DContMorphism.toCofunctor {C D : DirectedContainer}
    (phi : DContMorphism C D) : Cofunctor C D where
  obj := phi.onShapes;  lift := phi.onPos
  c0 := phi.m0;  c1 := phi.m1;  c2 := phi.m2

def Cofunctor.toDContMorphism {C D : DirectedContainer}
    (phi : Cofunctor C D) : DContMorphism C D where
  onShapes := phi.obj;  onPos := phi.lift
  m0 := phi.c0;  m1 := phi.c1;  m2 := phi.c2

theorem DContMorphism.toDContMorphism_toCofunctor {C D : DirectedContainer}
    (phi : DContMorphism C D) : phi.toCofunctor.toDContMorphism = phi

theorem Cofunctor.toCofunctor_toDContMorphism {C D : DirectedContainer}
    (phi : Cofunctor C D) : phi.toDContMorphism.toCofunctor = phi
\end{lstlisting}

\paragraph{Haskell.}
\begin{lstlisting}[style=haskell]
-- Both conversions are identity functions on the underlying data.
toCofunctor       :: DContMorph c d -> Cofunctor c d
toCofunctor m = Cofunctor (onShapes m) (onPos m)

toDContMorphism   :: Cofunctor c d -> DContMorph c d
toDContMorphism f = DContMorph (obj f) (lift f)
\end{lstlisting}


%===================================================================
\section{The category of directed containers (via cofunctors)}
%===================================================================

% --- 11. Cofunctor.id ----------------------------------------------
\subsection{Definition: \lean{Cofunctor.id}}

\begin{defn}[Identity cofunctor]
The identity cofunctor on $C$: object map is $\mathrm{id}$, lift is $\mathrm{id}$.
All laws hold by $\mathrm{rfl}$.
\end{defn}

\paragraph{Lean.}
\begin{lstlisting}[style=lean]
def Cofunctor.id (C : DirectedContainer) : Cofunctor C C where
  obj s := s
  lift _ p := p
  c0 _ _ := rfl;  c1 _ := rfl;  c2 _ _ _ := rfl
\end{lstlisting}

\paragraph{Haskell.}
\begin{lstlisting}[style=haskell]
idCofunctor :: Cofunctor c c
idCofunctor = Cofunctor id (\_ p -> p)
\end{lstlisting}


% --- 12. Cofunctor.comp_transport ----------------------------------
\subsection{Lemma: \lean{Cofunctor.comp\_transport}}

\begin{prop}[Transport commutes with composition]
Let $\psi : \mathcal{C}(D) \to \mathcal{C}(E)$.  Given shape equalities
$e : y = x$ in $D$ and $e_2 : \psi.\mathrm{obj}(x) = z$ in $E$, and a
position $k : P_E(z)$:
\[
  e^{-1} \triangleright \psi.\mathrm{lift}_x(e_2^{-1} \triangleright k)
  \;=\;
  \psi.\mathrm{lift}_y((\mathrm{congrArg}(\psi.\mathrm{obj}, e) \cdot e_2)^{-1} \triangleright k).
\]
This is the key dependent step enabling composition of cofunctors.
\end{prop}

\paragraph{Lean.}
\begin{lstlisting}[style=lean]
theorem Cofunctor.comp_transport {D E : DirectedContainer} (psi : Cofunctor D E)
    {x y : D.Shape} {z : E.Shape} (e : y = x) (e2 : psi.obj x = z)
    (k : E.Pos z) :
    e.symm [> psi.lift x (e2.symm [> k)
      = psi.lift y (((congrArg psi.obj e).trans e2).symm [> k)
\end{lstlisting}

\paragraph{Haskell.}
\begin{lstlisting}[style=haskell]
-- Transport-coherence lemma; collapses to rfl in Lean.
-- No direct Haskell analogue (transports are trivial without dependent types).
\end{lstlisting}


% --- 13. Cofunctor.comp -------------------------------------------
\subsection{Definition: \lean{Cofunctor.comp}}

\begin{defn}[Composition of cofunctors]
Given cofunctors $\varphi : C \to D$ and $\psi : D \to E$, their composite
$\psi \circ \varphi : C \to E$ has:
\begin{align*}
  \mathrm{obj}(s) &= \psi.\mathrm{obj}(\varphi.\mathrm{obj}(s)), \\
  \mathrm{lift}_s(m) &= \varphi.\mathrm{lift}_s(\psi.\mathrm{lift}_{\varphi.\mathrm{obj}(s)}(m)).
\end{align*}
The lift composes \emph{backwards} (positions pulled through $\psi$ then $\varphi$).
\end{defn}

\paragraph{Lean.}
\begin{lstlisting}[style=lean]
def Cofunctor.comp {C D E : DirectedContainer}
    (phi : Cofunctor C D) (psi : Cofunctor D E) : Cofunctor C E where
  obj s := psi.obj (phi.obj s)
  lift s m := phi.lift s (psi.lift (phi.obj s) m)
  c0 s m := (congrArg psi.obj (phi.c0 s (psi.lift (phi.obj s) m))).trans
    (psi.c0 (phi.obj s) m)
  c1 s := by rw [psi.c1, phi.c1]
  c2 s h k := by
    rw [psi.c2, phi.c2,
      psi.comp_transport (phi.c0 s (psi.lift (phi.obj s) h)) (psi.c0 (phi.obj s) h)]
\end{lstlisting}

\paragraph{Haskell.}
\begin{lstlisting}[style=haskell]
compCofunctor :: Cofunctor c d -> Cofunctor d e -> Cofunctor c e
compCofunctor phi psi = Cofunctor
  { obj  = obj psi . obj phi
  , lift = \s m -> lift phi s (lift psi (obj phi s) m)
  }
\end{lstlisting}


% --- 14--16. Cofunctor category laws ------------------------------
\subsection{Theorems: category laws for \lean{Cofunctor}}

\begin{thm}[Left unit] $\mathrm{comp}(\mathrm{id}_C, \varphi) = \varphi$ (by \lean{rfl}).
\end{thm}
\begin{thm}[Right unit] $\mathrm{comp}(\varphi, \mathrm{id}_D) = \varphi$ (by \lean{rfl}).
\end{thm}
\begin{thm}[Associativity]
$\mathrm{comp}(\mathrm{comp}(\varphi, \psi), \chi) = \mathrm{comp}(\varphi, \mathrm{comp}(\psi, \chi))$
(by \lean{rfl}).
\end{thm}

\paragraph{Lean.}
\begin{lstlisting}[style=lean]
theorem Cofunctor.id_comp {C D : DirectedContainer} (phi : Cofunctor C D) :
    comp (id C) phi = phi

theorem Cofunctor.comp_id {C D : DirectedContainer} (phi : Cofunctor C D) :
    comp phi (id D) = phi

theorem Cofunctor.comp_assoc {B C D E : DirectedContainer}
    (phi : Cofunctor B C) (psi : Cofunctor C D) (chi : Cofunctor D E) :
    comp (comp phi psi) chi = comp phi (comp psi chi)
\end{lstlisting}

\paragraph{Haskell.}
\begin{lstlisting}[style=haskell]
-- prop_id_comp phi = compCofunctor idCofunctor phi == phi
-- prop_comp_id phi = compCofunctor phi idCofunctor == phi
-- prop_assoc phi psi chi =
--   compCofunctor (compCofunctor phi psi) chi
--     == compCofunctor phi (compCofunctor psi chi)
\end{lstlisting}


%===================================================================
\section{Directed-container morphisms as a category}
%===================================================================

% --- 17. DContMorphism.id / comp / functoriality -----------------
\subsection{Definitions and theorems: \lean{DContMorphism} as a category}

\begin{defn}[Identity directed-container morphism]
$\mathrm{id}_C = (\mathrm{Cofunctor.id}\;C).\mathrm{toDContMorphism}$ --- the identity on data.
\end{defn}

\begin{defn}[Composition of directed-container morphisms]
$\mathrm{comp}(\varphi, \psi) = (\varphi.\mathrm{toCofunctor} \circ \psi.\mathrm{toCofunctor}).\mathrm{toDContMorphism}$.
\end{defn}

\begin{thm}[\lean{toCofunctor} preserves identities]
$(\mathrm{id}\;C).\mathrm{toCofunctor} = \mathrm{Cofunctor.id}\;C$ (by \lean{rfl}).
\end{thm}

\begin{thm}[\lean{toCofunctor} preserves composition]
$(\varphi \cdot \psi).\mathrm{toCofunctor} = \varphi.\mathrm{toCofunctor} \circ \psi.\mathrm{toCofunctor}$
(by \lean{rfl}).

Together with the \lean{rfl} bijection, this makes \lean{toCofunctor} an
isomorphism of categories between directed-container morphisms and cofunctors.
\end{thm}

\paragraph{Lean.}
\begin{lstlisting}[style=lean]
def DContMorphism.id (C : DirectedContainer) : DContMorphism C C :=
  (Cofunctor.id C).toDContMorphism

def DContMorphism.comp {C D E : DirectedContainer}
    (phi : DContMorphism C D) (psi : DContMorphism D E) : DContMorphism C E :=
  (phi.toCofunctor.comp psi.toCofunctor).toDContMorphism

theorem DContMorphism.toCofunctor_id (C : DirectedContainer) :
    (id C).toCofunctor = Cofunctor.id C

theorem DContMorphism.toCofunctor_comp {C D E : DirectedContainer}
    (phi : DContMorphism C D) (psi : DContMorphism D E) :
    (phi.comp psi).toCofunctor = phi.toCofunctor.comp psi.toCofunctor
\end{lstlisting}

\paragraph{Haskell.}
\begin{lstlisting}[style=haskell]
idDCM :: DContMorph c c
idDCM = toDContMorphism idCofunctor

compDCM :: DContMorph c d -> DContMorph d e -> DContMorph c e
compDCM phi psi = toDContMorphism (compCofunctor (toCofunctor phi) (toCofunctor psi))
\end{lstlisting}


%===================================================================
\section{Translation conventions}
%===================================================================

\begin{enumerate}
\item \textbf{Dependent position maps.}  M2 and C2 involve transporting a
  position along a shape equality (M0/C0) before it can be consumed by shift.
  Haskell renders this as a plain function application; the type-level
  correctness guarantee is lost.

\item \textbf{Identity bijection.}  The isomorphism between
  \lean{DContMorphism} and \lean{Cofunctor} is witnessed by \lean{rfl} in both
  directions.  This means the two Lean structures carry bit-for-bit identical
  data; the distinction is purely conceptual (``container morphism'' vs.\
  ``cofunctor'').

\item \textbf{Round-trip gap.}  The equivalence $\mathrm{DCont} \simeq \mathrm{Cat}$
  (Ahman--Uustalu 2016, Thm~3.6) holds only up to bijection on position types,
  not definitional equality.  The file exhibits the bijection explicitly rather
  than claiming a \lean{rfl}-equality that does not hold.

\item \textbf{Proof bodies omitted.}  Only statements are translated.
\end{enumerate}

\end{document}
